\section{Introduction}\label{sec:introduction}

For a real number $t$, write $\|t\|$ for its distance to the nearest
integer, and let $\phi$ denote Euler's totient function. We study
approximation with a fixed additive shift and an arbitrary sequence of
tolerances. The exceptional null set in the result below is allowed to
depend on both of these data.

\begin{theorem}[Weak inhomogeneous Duffin--Schaeffer]
\label{thm:main}
Let $\gamma\in\mathbb R$ be fixed, and let
$\psi:\mathbb N\to[0,\infty)$ be finite-valued. If
\begin{equation}\label{eq:main-divergence}
  \sum_{q=1}^{\infty}\frac{\phi(q)}q\,\psi(q)=\infty,
\end{equation}
then, for Lebesgue-almost every $x\in\mathbb R$, the inequality
\[
  \|qx-\gamma\|<\psi(q)
\]
holds for infinitely many positive integers $q$.
\end{theorem}

Theorem~\ref{thm:main} resolves the weak inhomogeneous Duffin--Schaeffer
conjecture positively, including every prescribed irrational shift.
The integer nearest to $qx-\gamma$ is unrestricted: no coprimality
condition is imposed on the numerator. The conclusion is a divergence
statement, not a convergence converse or a quantitative asymptotic.

The mass transference principle of Beresnevich and Velani gives the
following Hausdorff-measure consequence.

\begin{corollary}[Hausdorff-measure divergence]\label{cor:hausdorff}
Fix $\gamma\in\mathbb R$ and a finite-valued
$\psi:\mathbb N\to[0,\infty)$. Let
$f:(0,\infty)\to(0,\infty)$ be continuous and nondecreasing, with
$f(r)\to0$ as $r\downarrow0$, and suppose that $f(r)/r$ is monotone.
Set $f(0)=0$ and
\[
  W_\gamma(\psi)=\{x\in\mathbb R:
    \|qx-\gamma\|<\psi(q)\text{ for infinitely many }q\}.
\]
If
\[
  \sum_{q=1}^{\infty}\phi(q)f\!\left(\frac{\psi(q)}q\right)=\infty,
\]
then, for every bounded interval $I$,
\[
  \mathcal H^f(I\cap W_\gamma(\psi))=\mathcal H^f(I),
\]
where $\mathcal H^f$ denotes Hausdorff $f$-measure.
\end{corollary}

The corollary is a sufficient divergence criterion for each fixed
shift, with unrestricted numerators. It requires no monotonicity of
$\psi$ and asserts neither a convergence converse nor a coprime-numerator
version. Its proof is given in Section~\ref{sec:consequences}.

\subsection{Background and prior work}

The starting point is the monotone theory of metric approximation.
Khintchine's theorem~\cite{Khintchine1924}, in its usual modern form,
gives a full-measure conclusion for $\|qx\|<\psi(q)$ when $\psi$ is
non-increasing and $\sum_q\psi(q)$ diverges. Szüsz extended the
monotone theory to each fixed inhomogeneous shift~\cite{Szusz1958};
see also the precise modern formulation in~\cite{Yu2021}.
Removing monotonicity changes the problem because many denominators
can describe strongly overlapping approximation intervals.
Duffin and Schaeffer~\cite{DS1941} exhibited this obstruction in the
homogeneous setting and proposed using reduced fractions with the
weighted sum $\sum_q\phi(q)\psi(q)/q$. Koukoulopoulos and
Maynard~\cite{KM2020} proved their conjecture. The weight $\phi(q)/q$
records the proportion of reduced residue classes; it also gives a
meaningful sufficient divergence criterion when all numerators are
allowed.

In the inhomogeneous problem, even unrestricted numerators do not
make unweighted divergence sufficient. Ramírez~\cite[Theorem~1]{Ramirez2017}
constructed nonmonotone counterexamples for every prescribed countable
collection of shifts. His constructions have finite totient-weighted
sum, as noted in that paper, and thus leave the weighted question open.
Chow, Hauke, Pollington and
Ramírez~\cite[Theorem~1.5]{CHPR2025} show that, for suitable shifts
including all rational shifts, even an approximation function that is
non-increasing on its support can have divergent unweighted sum and a
null approximation limsup. These results explain why a proof must
control arithmetic overlap on sparse supports, not merely sum the
lengths of the intervals.

The fixed-shift weighted question appears explicitly in
Yu~\cite[Question~1.2]{Yu2021} and under the name \emph{weak inhomogeneous
Duffin--Schaeffer conjecture} in Chow and
Technau~\cite[Conjecture~1.22]{ChowTechnau2024}. It is also
Conjecture~2 of Beresnevich, Hauke and Velani~\cite{BHV2024}.
The word ``weak'' concerns the conclusion: the integer $a$ in
$|qx-a-\gamma|<\psi(q)$ is arbitrary. The stronger formulation
requires $(a,q)=1$; this distinction is already explicit in
Ramírez~\cite{Ramirez2017}. Theorem~\ref{thm:main} addresses the weak
formulation for every prescribed real shift.

The coprime-numerator assertion is false in general: Hauke-Treuer,
Maynard and Pollington~\cite[Theorem~1]{HMP2026} construct
counterexamples for every nonzero rational shift and certain Liouville
shifts. Independent work of He and Liao~\cite[Theorems~1--2]{HL2026}
gives counterexamples for every nonzero rational shift and for a
residual set of real shifts.

Several earlier advances treat important parts of the nonmonotone
problem. Yu's Fourier approach~\cite{Yu2019} gives inhomogeneous
results under additional divergence assumptions, and his
Erdős--Vaaler type theorem~\cite[Theorem~1.3]{Yu2021} treats
$\psi(q)=O(1/(q(\log\log q)^2))$, with divergent unweighted sum,
for a class of irrational shifts containing all non-Liouville
irrationals. Chow and Technau~\cite[Theorem~1.23]{ChowTechnau2024}
prove the weak conjecture for certain structured nonmonotone
functions arising in multiplicative approximation on fibres.
Those theorems place restrictions on the shift, the approximation
function, or both; they do not supply the arbitrary-function,
arbitrary-fixed-shift assertion used here.

Beresnevich, Hauke and Velani~\cite[Theorem~14]{BHV2024} establish
the weak conjecture for every rational shift. More precisely, for
$\gamma=A/B$ in lowest terms with $B>0$, their solutions can be
required to satisfy $(A+aB,q)=1$. This is a restriction adapted to
the shift and is different from ordinary numerator coprimality
$(a,q)=1$. Their work also gives irrational-shift results under
additional block or approximation hypotheses, and constructs
classes of suitable shifts. The remaining task addressed by
Theorem~\ref{thm:main} is to retain a completely arbitrary fixed
irrational shift together with arbitrarily sparse, nonmonotone
approximation functions.

\subsection{Ancestry of the proof methods}

The arithmetic part adapts the common-pivot approach to prime
valuations developed in the homogeneous Duffin--Schaeffer work of
Koukoulopoulos and Maynard~\cite{KM2020}, simplified by Green and
Walker~\cite[Proposition~1.2 and Section~3]{GW2021}, and used in the
weighted argument of Hauke-Treuer, Vazquez and
Walker~\cite[Proposition~2.1]{HVW2026}. In that approach, failure of an
arithmetic estimate concentrates the exponents of two denominator
families around one common integer. The cited common-pivot
propositions concern hypothetical minimal counterexamples in their
respective arguments. We prove our weighted lcm estimate locally,
with its own totient capacities and multiplicative inflation
hypotheses, and then prove the small-scale extraction needed for
inhomogeneous overlaps. The one-coordinate concentration argument
is credited again at its point of use in Section~\ref{sec:preliminaries}.

Two further precedents help explain the organization. Chow and
Technau's shift-reduced numerator restrictions and localized full-measure
criteria~\cite[Definition~1.20 and Section~2.4]{ChowTechnau2024}
illustrate how numerator conditions can be adapted to inhomogeneous
arithmetic. Our conditions use rational models of the projected
phases and are specified independently below. Beresnevich, Hauke and
Velani obtain full measure by combining uniform distribution of selected
points $a/q$ with divergence or block-mass conditions and overlap
bounds~\cite[Theorems~11--12]{BHV2024}. Here the
transfer is proved directly for weighted finite arrays, using a
virtual comparison sum and a fixed finite-union approximation to an
avoided set. These are methodological comparisons; all the
construction and transfer estimates used in this paper are proved
here.

\subsection{Proof strategy and reusable estimates}

Put $r(q)=\psi(q)/q$. A row is the periodic array of points
$(z+\gamma)/q$, $z\in\mathbb Z$, with spatial half-width $r(q)$.
The proof works on arbitrarily late finite collections of rows whose
totient-weighted mass lies in a fixed small interval. On each collection
we construct a nonnegative weighted sum of interval indicators supported
inside the original approximation intervals. The construction changes
weights through prescribed prime-step updates and additional deletions.
For comparison, we omit only these additional deletions, keeping the
same realized prime-step update factors. This comparison sum is
equidistributed on fixed spatial
intervals; the expected discarded mass is small, and the retained sum
has a uniform second-moment bound. These bounds contradict the existence
of a positive-measure set avoiding all sufficiently late rows: approximate
that set by one fixed finite union of intervals, then use the second
moment to control the approximation error. Proposition~\ref{tab:conditional-reduction}
gives this transfer after the construction is defined.

There are two different obstructions to the second-moment bound. First,
nearness of two shifted points constrains an integer combination of
their numerators to an interval too short for direct averaging. The
inverse estimate of Proposition~\ref{ar:extraction} extracts denominators
concentrated near a common integer in prime-exponent coordinates.
Varying permitted factors then either bounds the overlap weight by a
rotation estimate or forces precise rational approximations to the grid
offsets. Section~\ref{sec:direct} treats both alternatives. Exact modeled
coincidences have total weight bounded linearly by the block's
totient-weighted mass, using explicit restrictions on numerator gcds.

Second, large common primes can create correlated alignments repeatedly
as different collections of rows begin to interact. Each row is assigned
an integer multiple of its denominator, called its centre. We expose the
prime factors of centres in an increasing sequence of scales. Before all factors have
been exposed, the row is represented by a finer projected grid containing
its eventual true points; its labels retain the original row and the
unexposed divisor data. At each prime step, short spatial cells group
projected points of the same width. Some entries are retained only when
their numerator is divisible by that prime. We pair portions of their
weights with portions of other entries in the cell; when a divisibility
test succeeds, the linked portion of the other weight is suppressed.
The precise rule in Section~\ref{sec:tables} preserves the first-moment
normalization while cancelling the most tightly correlated transfers.

The remaining correlations are paid by a squared-mass accounting
argument. Labels acquire equivalence classes through actual close
alignments, and later partitions preserve those classes even when
intermediate labels disappear. Deterministic interval grids record the
old connecting segments, allowing new short-cell partitions to respect the
history. At the first interaction of two previously separate collections,
a backward comparison transfers covariance terms to linked points that
passed their divisibility tests. The total incoming matching weight is
bounded by the weight available at each such point, as proved in
Lemma~\ref{join:capacity}. The averages use common residue translations;
they do not assume independent test outcomes or randomize the shift.

The reusable estimates are stated separately: small-scale extraction
with auxiliary weights in Proposition~\ref{ar:extraction} and
Corollary~\ref{ar:box-sampling}, the linear bound for exact rational-model
coincidences in Lemma~\ref{dir:exact-equalities}, and the variance budget
for history-preserving projected arrays in Lemma~\ref{tab:variance-budget}.

\subsection{Organization}

Section~\ref{sec:preliminaries} gives the finite-block reduction and
the elementary tensor and residue estimates used to select centres.
Sections~\ref{sec:centres}
and~\ref{sec:prescriptions} assign centres and prescribe the initial
weights. Section~\ref{sec:tables} gives the projected-array construction
and reduces the theorem to three precisely stated comparison and
partition assertions. Section~\ref{sec:arithmetic} develops the
sieve, rotation, and small-scale concentration estimates for close pairs.
Section~\ref{sec:direct} applies them to the comparisons
for adjacent simultaneous insertions. Section~\ref{sec:sparse} controls alignments
during a prime step, and Section~\ref{sec:partitions} constructs the
required partitions. Section~\ref{sec:joins} proves the remaining
first-interaction estimates and completes the finite induction.
Section~\ref{sec:consequences} proves the Hausdorff-measure corollary.
